\clearpage\thispagestyle{empty}
\pdfbookmark[0]{Quellen zu Über die Grenzen der geozentrischen Positionen der Planeten (1804)}{grenzen1804-sources}
\begin{center}\textbf{Quellen zu Über die Grenzen der geozentrischen Positionen der Planeten (1804)}\end{center}
\noindent Diese KI-Übertragung in heutiges Deutsch folgt Carl Friedrich Gauß’ deutschem Original in den \textit{Werken}, Band~VI (Göttingen, 1874), S.~106--118, einschließlich der drei Tabellen für Pallas und Ceres. Der Aufsatz erschien zuerst als Beitrag XIV im Augustheft 1804 der von Franz Xaver von Zach herausgegebenen \textit{Monatlichen Correspondenz zur Beförderung der Erd- und Himmels-Kunde}, Band~10, S.~173--191. \href{https://gdz.sub.uni-goettingen.de/id/PPN236007467}{Historischer deutscher Text in den Werken}. \href{https://www.digitale-sammlungen.de/en/view/bsb10538602?page=181}{Erstdruck}.
\par\smallskip\noindent Gauß’ zwei Verweise auf seinen Aufsatz vom Mai 1804 beziehen sich auf S.~98 und 100 dieses Bandes. Im \href{https://www.digitale-sammlungen.de/en/view/bsb10538601?page=421}{Erstdruck vom Mai}, Band~9, Beitrag L, S.~385--400, entsprechen ihnen S.~385\,ff.\ und 394.
\par\smallskip\noindent Zur Terminologie der Bahnen siehe Gauß’ \href{https://www2.karlin.mff.cuni.cz/~tuma/Aplikace21/theoryofmotionof00gausuoft.pdf#page=28}{\textit{Theory of the Motion of the Heavenly Bodies} in Charles Henry Davis’ englischer Übersetzung} (Boston, 1857), Buch I, Abschnitt I, Art.~1--3 und den Anfang von Art.~4, und Ernst Friedrich Maria Jochims \href{https://elib.dlr.de/94028/1/SAB-III-gesamt-2022-02-09.pdf#page=74}{\textit{Satellitenbewegung}, Band III} (PDF vom 9.\ Februar 2022), S.~60--63.
\par\smallskip\noindent Übertragung und Satz wurden mit KI erstellt.
\clearpage
